% ==============================================================================
% SLIDE 1: ENVIRONMENTAL DETERMINISM & MATERIAL SCARCITY
% ==============================================================================
\begin{frame}[label=scarcity]{The Statics of Scarcity: Materials \& Form}
  \begin{columns}[T, onlytextwidth]
    % Left Column: Photo / Technical Diagram
    \begin{column}{0.48\textwidth}
      \slidephoto{Yposkafa are Cave homes carved into volcanic pumice, from paulshark and Getty Images.jpg}{Subterranean \textit{yposkafa} carved into volcanic pumice (\textit{aspa}) along the Oia cliff face}
    \end{column}

    % Right Column: Engineering Bullet Points
    \begin{column}{0.48\textwidth}
      \textbf{\color{aegeanblue}Material Constraints Dictate Geometry:}
      \vspace{0.2em}
      \begin{itemize}
          \setlength{\itemsep}{0.15em}
        \item \textbf{The Timber Deficit:} Total absence of structural wood forced the abandonment of post-and-beam frames and flat joist floors.
        \item \textbf{Pure Compression Geometry:} Stone barrel vaults and hemispherical domes distribute ceiling loads as pure compressive thrust.
        \item \textbf{Excavation vs. Erection:}
          \begin{itemize}
            \item \textit{Santorini (Yposkafa):} Carved into compacted volcanic pumice (\textit{aspa}); the earth itself acts as the structural arch.
            \item \textit{Granite Islands (Ktista):} Freestanding thick stone walls (60--80\,cm) stacked from metamorphic rock.
          \end{itemize}
        \item \textbf{Battered Profiles:} Thickened wall bases lower the center of gravity to resist lateral seismic shear.
      \end{itemize}
    \end{column}
  \end{columns}
\end{frame}
